\scanpage{7}
\mspage{5}

% pg 60
\uline{pg 60}\qquad Je n'ai pas confiance en ``il est immédiat'' de 6.4.

% left of pg 60
% slanted
\margin{\tikz[baseline=(c.base)]{\node[draw,ellipse,inner sep=1pt] (c) {\shortstack{c'est pourtant\\trivial.}};}}

% pg 61, 61 bis
\uline{pg 61, 61 bis} : Je n'ai pas lu.

% next paragraph
Il me semble qu'il doit être possible de définir, à isogénie près, des ``parties finies'' et ``parties toriques'', ayant le bon rang --

% left of preceding paragraph
% slanted
\margin{\tikz[baseline=(c.base)]{\node[draw,ellipse,inner sep=1pt] (c) {\shortstack{non, mon\\argument\\marche}};}}

% pg 64
\uline{pg 6\deleted{\unsure{3|5}}4/}\deleted{\unsure{6.6|66|68}} : \deleted{\unsure{$l$ 10|$l$ 11|Serre}} (réf. à Serre pour l'indép. de $l$ ligne 11 ?)

% pg 64, second line
Un peu malhonnête d'être si bref, car 6.6 est utilisé en 5.13.

% left of pg 64
% slanted
\margin{\tikz[baseline=(c.base)]{\node[draw,ellipse,inner sep=1pt] (c) {que faire ?};}}

% division
\hrule\smallskip\hrule\smallskip\hrule

% section 7, p 65
\tikz[baseline=(c.base)]{\node[draw,ellipse,inner sep=1pt] (c) {7};}\quad \uline{p 65}.\qquad $A_0$ n'est pas défini\qquad \margin{si}

% p. 67
\margin{\tikz[baseline=(c.base)]{\node[draw,ellipse,inner sep=1pt] (c) {?};}}\uline{p. 67}\qquad (7.1.4)\qquad français --

% division
\hrule\smallskip\hrule\smallskip\hrule

% section 8
\tikz[baseline=(c.base)]{\node[draw,ellipse,inner sep=1pt] (c) {8};}

% division
\hrule\smallskip\hrule\smallskip\hrule

% section 9, p 80
\tikz[baseline=(c.base)]{\node[draw,ellipse,inner sep=1pt] (c) {9};}\quad \uline{p 80}\qquad 9.2.3\qquad $u$ plutôt que $u(g)$

% p 82
\uline{p 82}\qquad (9.3.3)\qquad $+$ compatibilité avec $R$ ? --

% p 87
\uline{p 87}\qquad 9.4.\qquad $S$ \uuline{strictement} hensélien ? (on le suppose en p 88, $l-9$)

% next indented line
pas d'hyp. d'excellence ?

% p 91
\uline{p 91} :\quad 9.5.8.\qquad Ne mettrait-tu pas une génuflexion ``au signe près''.

% pg 92 bis
\uline{pg 92 bis}\qquad $S$ est-il strictement hensélien ?

% division
\hrule\smallskip\hrule\smallskip\hrule

% section 10, pg 98
\tikz[baseline=(c.base)]{\node[draw,ellipse,inner sep=1pt] (c) {10};}\quad \uline{pg 98}\qquad 10.5.3.\quad Une extension de $K$ (qui ici est infini) est-elle requise ?? .

% last line
10.5.4.\qquad Dire \uline{où} on fera ces vérifications --
% archive sheet number: 6
